\documentclass[11pt,a4paper]{article}
% Compile with XeLaTeX (fontspec and xeCJK).
\usepackage{fontspec}
\setmainfont{LibertinusSerif}[Extension=.otf,UprightFont=*-Regular,BoldFont=*-Bold,ItalicFont=*-Italic,BoldItalicFont=*-BoldItalic]
\setsansfont{LibertinusSans}[Extension=.otf,UprightFont=*-Regular,BoldFont=*-Bold,ItalicFont=*-Italic]
\usepackage{xeCJK}
\setCJKmainfont{Hiragino Mincho ProN}
\xeCJKDeclareCharClass{Default}{"00B7, "2014, "2018 -> "201F, "2026}
\usepackage[english]{babel}
\usepackage[a4paper,margin=2.5cm,top=2.3cm]{geometry}
\usepackage{graphicx}
\usepackage{xcolor}
\usepackage{titlesec}
\usepackage{fancyhdr}
\usepackage[hidelinks]{hyperref}
\definecolor{accent}{RGB}{31,73,125}
\definecolor{muted}{RGB}{93,104,119}
\hypersetup{colorlinks=true,urlcolor=accent,linkcolor=accent,pdfauthor={Romain Fernex},pdftitle={The origins of Korea and the kingdom of Goguryeo}}
\providecommand{\tightlist}{\setlength{\itemsep}{0pt}\setlength{\parskip}{0pt}}
\providecommand{\ul}[1]{\underline{#1}}
\setcounter{secnumdepth}{0}
\titleformat{\section}{\large\bfseries\color{accent}}{}{0em}{}
\setlength{\parindent}{1.2em}
\setlength{\parskip}{0.35em}
\pagestyle{fancy}
\fancyhf{}
\fancyhead[L]{\small\color{muted}Romain Fernex}
\fancyhead[R]{\small\color{muted}\itshape The origins of Korea and the kingdom of Goguryeo}
\fancyfoot[C]{\small\color{muted}\thepage}
\renewcommand{\headrulewidth}{0pt}
\newcommand{\articlefigure}[2]{%
  \begin{center}
  \includegraphics[width=0.55\linewidth,height=6.5cm,keepaspectratio]{#1}%
  \ifx\relax#2\relax\else\par\smallskip{\small\itshape\color{muted}#2\par}\fi
  \end{center}}

\begin{document}
\thispagestyle{empty}
\noindent{\small\sffamily\color{accent}\MakeUppercase{History of Korea}}\par\medskip
\noindent{\LARGE\bfseries The origins of Korea and the kingdom of Goguryeo\par}
\smallskip\noindent{\large\itshape Between nationalist narratives and historical realities\par}
\medskip
\noindent{\color{accent}\rule{\linewidth}{0.8pt}}\par\smallskip
\noindent{\small\color{muted}Romain Fernex · 2022 · Article in English}\par\bigskip

``The strongest cause for feelings of nationality ... is identity of political antecedents; the possession of national history, and consequent community of recollections; collective pride and humiliation, pleasure and regret, connected with the same incidents in the past'', John Stuart Mill

In analysing the roots of national feeling in this way, Stuart Mill reminds us of a major, yet sometimes forgotten, role of history: its political role. Indeed, any nation, whether recently formed or not, that seeks to anchor itself in the long term often draws on history, and particularly on its instrumentalisation, as a way of asserting its right, its legitimacy, to exist or develop within a particular geographical or cultural space. The construction of an official history promoted by a government, sometimes far removed from, if not contrary to, established historical realities, thus very often supports claims within the nation itself as well as in its external relations.

In 2004, following applications by China and then North Korea, the archaeological sites of the tombs and capital of the ancient kingdom of Goguryeo{[}1{]}, located on either side of the border between North Korea and China, were inscribed on UNESCO's World Heritage List.

\articlefigure{../goguryeo-1/media/media/image1.jpeg}{Tomb of Gwanggaeto the Great, also known as the General's Tomb, dating from the 4th century BC}

This decision occasioned many controversies, concerning the application that led to it, the process that resulted in the designation, and the reactions of the parties involved after the decision was taken. Thus, the removal of all references to South Korea from Goguryeo's history as presented on the Chinese Ministry of Foreign Affairs website provoked both popular and diplomatic outrage among South Koreans, who felt that their history was being stolen. Beyond traditional diplomatic channels, numerous online forums saw Chinese and Korean participants clash, advancing arguments that were sometimes moderate but often steeped in strong nationalism and historical revisionism. On one side, Chinese internet users maintained that the kingdom of Goguryeo and the populations that preceded it were attached to China and therefore formed part of Chinese cultural unity; on the other, Korean users defended Goguryeo's role as a pillar of Korean identity and as one of the first concrete manifestations of Korea's control over the peninsula.

But what exactly was this kingdom of Goguryeo, and why does it arouse such passions? Throughout this article, we will try to offer some answers by examining, first, Goguryeo's origins, which occupy a particularly important place in this debate, especially on the Chinese side, and, second, the kingdom itself and its importance, not only in the region's historiography but also in claims to a Korean identity stretching back several millennia.

\section{The peninsula before Goguryeo}

400,000 BC: this is the date of the earliest traces of human populations on the Korean peninsula during the Palaeolithic. These were nomadic peoples, mainly from Manchuria and Northeast Asia, living primarily by hunting and fishing. The Jeulmun period began in the 6th millennium BC, marked by the appearance of comb-pattern pottery found mainly along the Han River, which flows through present-day Seoul.

\articlefigure{../goguryeo-1/media/media/image2.jpeg}{Pottery from the Jeulmun period}

Rice cultivation developed from 4000 BC, aided by a near-tropical summer climate accompanied by heavy monsoons in the same season, allowing some of these peoples to settle and their populations to grow. This development continued until 2000 BC with the transition to the Mumun period, itself marked by a new style of pottery with double-rimmed vessels. This period brought major discoveries conducive to settled life and the emergence of complex, organised civilisations, with the use of bronze in 700 BC and iron in 300 BC. China's growing influence in the region was also characteristic of this period, and it would play an increasing role over the following centuries. It was in this context that the kingdom of Gojoseon emerged in the northern peninsula, itself playing an important role in Korean identity and mythology.

Indeed, for Koreans, Gojoseon's origins lie in the myth of Dangun Waggeon, which dates its creation to 2333 BC, well before the beginning of the Mumun period. According to this myth, Dangun is Hwanung, son of Hwanin, the sky god in Cheondoism and Jeungsanism, two major religious currents in Korea besides Buddhism, which remains the dominant religion and identifies him with Indra, and of Ungnyo, a bear transformed into a woman by Hwanung. He founded the city of Asadal, Gojoseon's capital, on the site of present-day Pyongyang. Although most historians agree that Gojoseon was more akin to a grouping of tribes loosely connected to one another, it nevertheless holds particular importance in Korean culture. This first entity was conquered by Han China, ruled by Emperor Wu, in the 2nd century BC, and then came under its direct control for the first and last time in its history. The Han established four commanderies as administrative divisions to govern Gojoseon's territory.

\articlefigure{../goguryeo-1/media/media/image3.jpeg}{The four commanderies during the period of Chinese rule}

The commanderies of Lelang and Xuantu were the most important, with Lelang producing textiles and ceramics, skills that Koreans would preserve and develop throughout their history. This was therefore an important period of Sinicisation in the region, giving rise to numerous connections and exchanges between China and the populations of the Korean peninsula. It also provides substantial support for Chinese claims to include Korea within the sphere of Chinese civilisation. However, while China undeniably played a major role in this region at the time, this overlooks the many populations, clans and tribes that remained there and with which China maintained relations that were sometimes hostile, sometimes friendly. China indeed pursued a policy of appeasement, also called Heqin, consisting of showering certain clans with gifts to win their favour.

In the south were three Han tribes: the Mahan, the most numerous and based on agriculture; the Chinhan, dependent particularly on silk production; and the Pyonhan. These would give rise to several kingdoms clearly distinct from China a few centuries later. Alongside these three tribes were the Wa, the first population associated with Japan, who nevertheless retained a significant presence in the southeastern Korean peninsula.

Finally, in the north, on Gojoseon's territory, there were a total of nine tribes. Among them were Buyeo, a close ally of China, and Goguryeo. The latter appeared in 37 BC according to the \emph{Samguk Sagi} {[}2{]} (History of the Three Kingdoms), founded by Dongmyeong in southern Manchuria.

\articlefigure{../goguryeo-1/media/media/image4.jpeg}{King Dongmyeong of Goguryeo (58 BC -- 19 BC)}

A particularly warlike tribal grouping because its territory lacked sufficient resources to sustain its growing population, it rapidly conquered other territories and organised itself into four major administrative divisions. It even maintained diplomatic relations with China from 75 to 12 BC, and its ruler proclaimed himself king (wang in Korean) to emphasise his desire to distinguish himself from China. Indeed, although China sought the favour of certain clans, it also subjected many of the region's native populations to severe exploitation, provoking frequent revolts. Goguryeo's meteoric rise therefore worried China all the more because it was gaining influence among other local tribes opposed to China.

The breakup of the Han Empire in AD 220 gave many powers, including Goguryeo, the opportunity to free themselves from Chinese influence and develop an initial sense of Korean identity, or rather identities, given the still relatively heterogeneous character of these different tribes, despite their many similarities. This rise was short-lived with the ascent of China's Wei dynasty in northern China, which, fearing Goguryeo and wishing to pacify the region, attacked it and took control of its capital in 244. However, Goguryeo's disappearance was only temporary, a prelude to a rebirth barely a century later that would establish it as one of the peninsula's most powerful kingdoms for several centuries.

\section{The formation and early development of Goguryeo}

The populations living alongside one another before Goguryeo’s emergence help show both the extent and the limits of Chinese influence on the peninsula. Yet restricting ourselves to an analysis of this area's ethnic diversity does not fully address the geopolitical issues raised by the controversy over the \emph{Northeast Project}. Indeed, China's aim is to deny the political reality of a kingdom structured around a culture distinct from Chinese culture, laying the foundations of what would ultimately become the Korea we know. Likewise, for both South and North Korea, the aim is to root their existence and identity in ancient history by making this first kingdom of Goguryeo a cornerstone of their development.

It is therefore important to examine the kingdom itself and what made it distinctive, geopolitically, culturally and economically.

Let us therefore return to 37 BC, the historical and mythological date of Goguryeo's foundation. Its founder is often known as Chumong, also called Dongmyeong, and is said to have originated in the kingdom of Buyeo, founded in the 2nd century BC in northern Manchuria. At its emergence, Goguryeo was more akin to a tribal confederation structured around five Bus{[}3{]}, which would soon become administrative divisions of the kingdom. These five Bus themselves originated in smaller tribal groupings called Nas, which still retained relative autonomy at the beginning of this period. As the Bus increased their power, one in particular, the Kyeru Bu, supplanted the others and would provide the kings of the dynasties ruling Goguryeo. This grouping was both the product and the cause of a policy of assimilation and conquest of the region's other chiefdoms and tribes, which would form its sphere of influence. Among the nomadic tribes that fell under Goguryeo's rule, several, including the Khitan and Malgal, retained substantial autonomy from the five Bus, although they remained subject to royal power. However, for several decades, the geographical and political unit within which Goguryeo's population lived and with which it interacted most was the eumnak. Eumnaks were communities of varying size, under village chiefs, through which the work and lives of the entire common population were organised.

\articlefigure{../goguryeo-2/media/media/image1.png}{Map showing the main political entities of the Korean peninsula in 37 BC}

(the four commanderies of Han China are in purple)

Beyond these administrative considerations, an initial social hierarchy can already be observed. At the top of the social ladder were the Daega. This term could refer to the lord heading the Kyeru Bu, that is, the king, to the chiefs of the other Bus, or to local chiefs. Further down, the Sogas formed a property-owning petty bourgeoisie, generally receiving a salary from the government, not participating in agricultural activities, and able to aspire to a certain degree of social advancement. It should be noted that they could also be demoted to the lower class if they became too heavily indebted or disobeyed the court's orders. Most people belonged to the Haho category (tenants in Chinese), so called because its members owned no land. This class included peasants, artisans, merchants and others serving nobles but considered to be employed by the central government{[}4{]}, or directly under the control of great houses. The Haho class was far from homogeneous and included Homin, wealthy individuals who generally derived their fortunes from trade or crafts. Lower down the social ladder were the Yongjak, peasants who hired out their services to a wealthy landowner in exchange for a share of the harvest or a plot to cultivate. Finally, at the very bottom were the Nobi, servants directly employed by large landowners. These were often prisoners of war or peasants unable to repay their debts.

Economically, the royal authorities quickly introduced an initial system of taxation. Nevertheless, its effectiveness remained fairly limited, since it often had to operate through clan chiefs{[}5{]} and took the form of a simple tribute paid to the central authority.

Thus, almost from its emergence, Goguryeo society appeared as a complex system, different from that prevailing in China at the time and from the one China would establish after its first conquest of Goguryeo in 245. However, the effectiveness of the central authority claimed by the Daegas, and particularly the royal house, remained an open question. This problem was all the more acute because the kingdom's successive conquests had enabled agricultural advances and an expansion of cultivable land, fuelling Haho peasants' desire for individual ownership, although their labour was vital to the kingdom's wealth and functioning. Faced with peasants' growing demand for autonomy, a genuine centralising movement initiated by royal power enabled Goguryeo to emerge as a nation-state. This movement also arose from the need to develop an ambitious foreign policy, particularly in response to Han China, whose Liaodong commandery had harassed Goguryeo since its foundation and had already pushed it back to the Yalu valley. Far from a rapid transformation, it was the result of a slow, centuries-long process that began with the kingdom's creation and experienced several reversals throughout its history.

But how did this centralising movement unfold, and what evidence attests to its existence so early in Goguryeo's life?

An initial important element was the establishment of the political system of principalities. This centred on the king, who concentrated most power and ruled over the entire kingdom. Immediately beneath him were the chiefs of the five Bus forming the principalities, themselves divided into sub-principalities or Che-Bu. This first group of principalities covered most of the territory. Although four of these great principalities were subordinate to royal power, they retained an important consultative role and could sometimes exert real influence over the king's decisions. They coexisted with a second group comprising, on the one hand, other tribes accountable to the principalities (and especially to the king) but not directly attached to them (subordinate groups, see diagram below), and, on the other, chiefdoms that had fallen under Goguryeo's influence while retaining considerable autonomy in managing their territory and forming marquisates {[}6{]}.

\articlefigure{../goguryeo-2/media/media/image2.png}{Diagram of Goguryeo's administrative organisation before AD 245}

From: Noh Taedon, \emph{Korea's Ancient Koguryo Kingdom: A Socio-political History}, Global Oriental, 2014, p.97

This transformation of Goguryeo's political system was accompanied by the gradual disappearance of eumnaks under pressure from Haho seeking property ownership and a government that saw community chiefs as an obstacle to its power, preventing it from exercising authority directly over peasants. Several instruments therefore developed to address this disconnection between royal power and the population, the source of its economic and military strength.

The first measure used by the monarchy to unify the population was the introduction of a cult of the king and his descendants based on the myth surrounding the birth of Chumong, Goguryeo's founder, who was said to be the son of the agricultural goddess Habaek Nyo. As in many monarchies, this was also a powerful means of establishing rulers' legitimacy, and it amply fulfilled that role until Goguryeo's fall.

The second was an initial reform of taxation, gradually shifting from a relatively rudimentary system of tribute demands to a genuine, systematic tax-collection system, making it a much better instrument of fiscal policy. However, this transformation extended over several centuries and occurred more substantially during the provincialisation of Goguryeo's political system.

Finally, in international relations, Goguryeo distinguished itself as a remarkable military power, particularly through its cavalry and its mastery of ironworking, inherited in part from Buyeo, which had made this its speciality. Foreign trade also developed with the commanderies of the Wei dynasty, but these relations remained marked by mistrust, were often limited to tributes paid by China to certain tribes making up Goguryeo, frequently to encourage internal divisions, and could hardly obscure the climate of tension between the two kingdoms. These tensions culminated

in 245, when the Wei besieged and defeated Goguryeo's capital in retaliation for its policy of alliances with Chinese dynasties rival to the Wei.

\articlefigure{../goguryeo-2/media/media/image3.png}{Map of East Asia in AD 262 showing the most influential Chinese dynasties}

The 3rd century thus ended in a dark age for Goguryeo, but one that had seen the foundations of a nation-state emerge, particularly politically and, admittedly less markedly, economically. We also observe the beginnings of cultural unity through the cult devoted to the royal family, although it remained at a primitive stage.

\section{Goguryeo’s recovery and the Three Kingdoms}

Goguryeo’s recovery after the Wei conquest took place in a different geopolitical context. As China weakened in the second half of the third century, new regional powers began to shape the peninsula. Historians call this the Three Kingdoms period, a name derived from two much later texts, the \emph{Samguk Sagi} (Chronicle of the Three Kingdoms) and the \emph{Samguk Yusa} (Memorabilia of the Three Kingdoms).

\articlefigure{../goguryeo-3/media/media/image1.png}{Map of the Korean peninsula in AD 375}

First, how could Goguryeo, decapitated after its capital fell to Wei China\textbf{{[}7{]}} in the mid-3rd century, rise from its ashes? The main reason for this resurgence of influence lies in the geopolitical context of the time, particularly the power vacuum left by internecine wars between dynasties seeking power in China. Indeed, the Chinese Jin dynasty, which succeeded the Wei dynasty in 266, faced a series of civil wars from 291 onwards, known to historians as the ``War of the Eight Princes''. As its name suggests, this saw several princes of the Jin dynasty fight to seize power in place of Emperor Hui, who had a severe developmental disorder and was therefore considered unfit to rule. Barely five years after the end of the War of the Eight Princes, Jin China collapsed in favour of numerous barbarian kingdoms that divided its territory. This began a dark age for China, corresponding to the Sixteen Kingdoms period and lasting until 439. Thus, from the mid-3rd century, Goguryeo troops that survived the fall of its capital took advantage of China's chaos to recapture East Okjo and then East Ye and re-establish its presence across much of its former territory. This reconquest continued in 303, accompanied by renewed expansionism through the capture of the Lelang and Daifang commanderies, long central to maintaining Chinese authority in the region.

\articlefigure{../goguryeo-3/media/media/image2.jpeg}{Map of China's Sixteen Kingdoms in AD 350}

But Goguryeo's rebirth did not take place alone. It was also at this time\textbf{{[}8{]}} that three other kingdoms emerged, two of which helped give the Three Kingdoms period its name. First came Baekje, born in the Mahan region of the southern peninsula and probably originating among Han populations to which it was culturally close. In the southeast, in the Kyeongju basin, the Kim royal lineage founded Silla, which would emerge victorious at the end of the Three Kingdoms period, although we will return to this later. Finally, modest in size but also relatively important to the regional balance, Kaya occupied the far south of the peninsula in the Naktong River valley. Its population included a significant proportion of Wa\textbf{{[}9{]}}, the founding ethnic group of Japan and particularly of the Yamato kingdom.

The early 4th century thus differed from preceding periods through the emergence of two major new regional powers, Baekje and Silla, and the great return of Goguryeo, finally escaping competition from China. It was during this period that Goguryeo emerged as the region's pre-eminent power, and this new prosperity was accompanied by major economic, cultural and political transformations. We must now analyse these to understand their extent and hence their relevance to the origins of this sense of Koreanness, at the heart of the debate we are examining.

Let us first look at Goguryeo's new political organisation after several decades in China's shadow. Indeed, this period of Chinese ``domination'' saw the end, or at least the weakening, of the system of principalities in favour of a provincial system with provincial officials relaying Chinese central power and, after the reconquest, royal power. The eumnaks, already disappearing, vanished completely, and direct control over the population finally passed to the royal family. The five Bus were thus reorganised into provinces\textbf{{[}10{]}}, each headed by a Yoksal (governor) exercising authority over several towns led by Cheoryogu (magistrates), assisted by Karadal (aldermen). At the bottom of the administrative hierarchy, mainly in rural areas, Rucheo acted as prefects. Finally, all these officials were accountable to the king, the sole holder of authoritative and coercive power.

This substantial centralisation of power enabled the completion of tax reform, replacing the primitive model of tribute demands with systematic and, above all, progressive taxation (to a certain extent). Royal power thus became the sole holder of taxation rights in place of lords and other large landowners. Although it occasionally granted certain high-ranking officials or deserving vassals a sigeup, a plot of land on which the lord could also collect taxes, it received the vast majority of tax revenue. This taxation targeted agricultural income first and was restricted to independent adult men. Faced with particularly significant processes of social diversification and stratification at the time, a genuine consideration of taxpayers' economic circumstances developed. The appearance of social assistance programmes, allowing certain struggling households to receive grain rations in exchange for a debt to the government, thus coincided with this awareness and aimed to provide better care for the kingdom's human capital.

Also domestically, the relocation of the capital to Pyongyang in 343, following repeated attacks by Former Yan China, enabled the expansion of a universal culture specific to Goguryeo despite numerous borrowings from China. This development began with the acceptance of Buddhism as a pillar of this pan-Goguryeo culture, but also with the systematisation of the Chumong myth, which definitively incorporated all Goguryeo monarchs into the kingdom's official and, above all, legendary history. This culture nevertheless had several variants: one can distinguish that developing in the north under Pyongyang's influence from that appearing in the south, still influenced by Goguryeo's former capital in Jian. In many regions, Taoism thus mingled with Buddhism and beliefs indigenous to this part of the peninsula. In this respect, studying the murals in many tombs of Goguryeo rulers and aristocrats, with which we began our discussion, allows us to appreciate the artistic movements that passed through Goguryeo. This evolution is evident both in the materials used to create these sumptuous frescoes and in the iconography dominating each period and region. Thus, during this first period, extending from the early 4th to the mid-6th century, many paintings depict everyday life or the four pioneering deities of Korean mythology. Several elements of Buddhist iconography, such as the lotus flower, also recur in the burial chambers of the period.

\articlefigure{../goguryeo-3/media/media/image3.jpeg}{Mural depicting a hunting scene,}

\emph{found in a tomb dating from the Goguryeo period}

Thus, describing Goguryeo at this time as a mere administrative division of China seems largely unwarranted, or at least particularly misleading.

Nevertheless, now that we have seen how Goguryeo differed from China and other regional powers in many respects, an important instrument of power remains to be analysed: its military strength. Indeed, Goguryeo rapidly emerged as Korea's almost unchallenged regional military power, and above all one native to the peninsula, for nearly a hundred years after centuries of Chinese influence. Let us now see how Goguryeo managed to place itself at the head of the regional geopolitical order.

From 373, under King Sosurim\textbf{{[}11{]}}, and after a sometimes tumultuous beginning with repeated Yan invasions, Goguryeo managed to pacify its relations with Qin China and could finally turn its dreams of expansion southward. Although it had already successfully attacked Silla in 345, obtaining an agreement in 348 in which the latter recognised Goguryeo's military superiority, it now sought to conquer Baekje, which had resisted it until then. This resurgence was marked by the successful capture in 375 of Sugok fortress, held by Baekje, which attempted a retaliatory campaign that failed completely. Yet the true turning point in Goguryeo's expansionist policy came in 391 with the accession of King Kwang-Gae-Toh, Goguryeo's most famous ruler after Chumong. He launched a wave of successful attacks on Baekje and, according to the Samguk Sagi, captured nearly 58 fortresses and 700 villages in less than a decade. At the same time, Goguryeo could count on the unwavering support of Silla, which appeared in many respects to be its vassal. A policy of hostage diplomacy\textbf{{[}12{]}} was also established between the two powers.

\articlefigure{../goguryeo-3/media/media/image4.jpeg}{Portrait of King Kwang-Gae-Toh}

This relationship of vassalage peaked in 400 when, during an attack on Silla by Japanese marauders led by Baekje and its faithful ally Kaya, Goguryeo came to Silla's aid and drove the enemy forces out of its territory. Following this event, Goguryeo troops remained stationed in Silla, and the hostage diplomacy already in place expanded significantly. However, tensions soon marked this relationship, particularly over the garrisons present in Silla, culminating in the accession of a king hostile to Goguryeo. Worse still, Silla drew closer to Baekje, and in the second half of the 5th century the two powers signed an agreement to counter Goguryeo's unstoppable rise.

In 450, an incident at Hasulla fortress, in which a Goguryeo soldier was killed by a Silla commander, revived these tensions, leading four years later to a pre-emptive Goguryeo attack on Silla. Events then followed in rapid succession, and in 464 the massacre of Goguryeo forces based in Silla's capital definitively ended any alliance between the two powers. Eleven years later, Goguryeo captured Baekje's capital, Hanseon, whose power declined with the outbreak of civil war and the assassination of its ruler, while also annexing part of Silla's territory. Goguryeo, now dominating the vast majority of the peninsula, was at the height of its power and negotiated as an equal with Northern Wei China.

But from the end of the 5th century, this dynamic began to reverse. Silla, which had previously intervened relatively little in the region's conflicts, including within its alliance with Baekje, adopted a much more interventionist policy and annexed most of Kaya in 532 to extend its influence. A little later, when Goguryeo attacked Baekje's Toksan fortress, Silla deployed its army and defeated it. Yet it was 551 that marked the definitive end of Goguryeo's golden age, with the capture of the Han River basin (the river flowing through present-day Seoul) by the Silla-Baekje alliance.

Romain Fernex, dual-degree student at ESCP and Korea University

\section{Notes}

\textbf{{[}1{]}} For consistency, I use throughout this article the Revised Romanisation of Korean, which has been Korea's official romanisation system since the 2000s and is also used by Kyung Moon Hwang. However, a substantial proportion, if not the majority, of the articles I used in writing this piece employ different spellings; for instance, Michael J. Seth uses the McCune-Reischauer system, which was authoritative in English-language academia when he wrote his book. Thus, Goguryeo may appear as Koguryo ̆ and, more generally, the names of places, kingdoms, key figures and so on are spelt differently from those in this article.

\textbf{{[}2{]}} The \emph{Samguk Sagi}, or History of the Three Kingdoms in Korean, is one of the major works, alongside the \emph{Samguk Yusa} (Memorabilia of the Three Kingdoms), published in 1279, dealing with the region's history in this period, reaching back several centuries before our era. Written in 1145 by the Korean historian and general Kim Busik, it is indeed one of the main written sources on the Korean peninsula available to historians of this period\ldots{} Yet certain passages are particularly romanticised and strongly influenced by the period in which the text was written, a period very different from the one it describes. Moreover, the foundational text of Japanese history, the \emph{Nihon Shoki}, and several major Chinese texts, including the \emph{Shiji} by the Chinese historian Sima Qian, himself a contemporary of this period, published in 100 BC, are also valuable sources for documenting it.

\textbf{{[}3{]}} Although the term Bu as an administrative division of Goguryeo is uncontroversial, the same cannot be said of the interpretation of its meaning or meanings in other contexts. When it designates the entity preceding these administrative divisions, it is particularly difficult to define. Some historians initially suggested that a Bu in its earliest form was a tribal grouping, but the term tribe can sometimes pose a problem. Indeed, following L. H. Morgan's definition in \emph{Ancient Society} in 1877, a tribe can be defined as an intermediate state between a classless, disorganised society and a nation-state. Yet a relatively developed hierarchy can be observed very early within the Bus, potentially calling this initial definition of a Bu as a tribal grouping into question.

\textbf{{[}4{]}} When I say that some Haho were employed by the central government, this clearly does not mean a direct relationship of exploitation, and it often operated through great local families. However, the major difference from the second category of Haho introduced immediately afterwards is that they served the interests of the state rather than private interests.

\textbf{{[}5{]}} This is all the more true of the royal authorities' relations with the nomadic tribes subject to them but retaining considerable autonomy, such as the Khitans or the Malgals. In these particular cases, the tribute collected was much smaller than for the eumnaks and always passed through the clan chiefs, who thus acted as intermediaries. However, this autonomy was offset by the fact that the state could mobilise them at will for military or economic purposes whenever the situation required it.

\textbf{{[}6{]}} The term is used by Noh Taedon in \emph{Korea's Ancient Koguryo Kingdom: A Socio-political History}, but he explains that, although it has no basis in the Korean literature of the time, it is the most appropriate term here, given the strong similarities in functioning between the ``marquisates'' mentioned here and their European counterparts, which appeared much later.

\textbf{{[}7{]}} This refers to China under the Cao Wei dynasty, lasting from AD 220 to 266. It should not be confused with China under the Northern Wei (386-535) or Western Wei (535-557) dynasties, which appeared much later and are also discussed further on in this article.

\textbf{{[}8{]}} Although the Samguk Sagi dates Silla's and Baekje's foundations to 57 BC and 18 BC respectively, no historical evidence currently corroborates this information, and most historians consider these kingdoms' emergence to be well before the date given for Goguryeo in that same work, which is itself relatively plausible. Thus, Silla's and Baekje's emergence was probably contemporary with Goguryeo's rebirth rather than its foundation.

\textbf{{[}9{]}} The Nihon Shoki also mentions a significant presence of these Wa peoples on the peninsula through the mythical kingdom of Minima, said to have belonged to Yamato and to have been lost to Silla. According to historians, however, this was largely imagined, and the existence of a fully fledged Wa kingdom on the peninsula is particularly difficult to prove given current knowledge of this period.

\textbf{{[}10{]}} Several competing theories exist about what exactly these new administrative provinces designated. For example, some suggest Goguryeo was divided between Naepyeong, the area directly influenced by the capital, and Waepyeong, covering the rest of the territory, itself divided into five provinces, each headed by a Yoksal. Others, however, consider that there were ten provinces in total, hence ten Yoksals, and so on\ldots{}

\textbf{{[}11{]}} It was also under this king that Buddhism was officially accepted in Goguryeo and the first legal code was published, laying the foundations of Goguryeo as it existed at its peak.

\textbf{{[}12{]}} Hostage diplomacy involves a country sending high-ranking young people, often members of a royal lineage, to an allied country as guarantees to prevent the alliance from being broken. This same policy would eventually work against Goguryeo by fuelling the division between Silla's pro-Goguryeo and anti-Goguryeo camps.

\section{Bibliography}

\begin{itemize}
\item
  Michael J.Seth, \emph{A history of Korea -- From Antiquity to the present}, 2010, Rowman \& Littlefield Publishers
\item
  Kyung Moon Hwang, \emph{A History of Korea}, 3rd edition, 2021, Palgrave
\item
  Kim Taesik, \emph{Koguryeo and Gaya - contacts and consequences}, Hungik universitt, journal of Northeast Asian History, Vol. 4 n°1 (p41-82), 2007
\item
  Jung Woon Yong, \emph{Trends in Koguryo's relationship with Paekche and Silla during the 4th -- 7th centuries}, International Journal of Korean history, Vol. 8, 2006
\item
  Mark Byington, \emph{The War of Words Between South Korea and China Over An Ancient Kingdom : Why Both sides are misguided}, history news network, Columbian college of arts \& science, The George Washington university, 2004
\item
  Thomas Chase, \emph{Nationalism and the Net} \emph{: online discussion of Goguryeo history in China and South Korea}, Monash university, China Information, 2011
\item
  John Stuart Mill, Considerations on representative government, Ch. XVI, Of Nationality, as connected with representative government, 1861
\item
  Noh Taedon, \emph{Korea's Ancient Koguryo Kingdom : A Socio-political History}, global oriental, 2014
\item
  Kim Chang-seok, \emph{Goguryeo Society and Its Economy}, The Review of Korean Studies, Vol.7 n°4, 2004
\item
  Jeon Ho-tae\emph{, Goguryeo Culture and Mural Paintings}, The Review of Korean Studies, Vol.7 n°4, 2004
\end{itemize}

\end{document}
